% Generated draft: requires mathematical and source review.
\documentclass[11pt]{article}
\usepackage[margin=25mm]{geometry}
\usepackage{fontspec}
\setmainfont{Latin Modern Roman}
\usepackage{amsmath,amssymb,mathtools,unicode-math}
\setmathfont{Latin Modern Math}
\usepackage{enumitem,needspace,xurl,hyperref}
\setlength{\parindent}{0pt}
\setlength{\parskip}{5pt}
\setlength{\emergencystretch}{3em}
\newcommand{\N}{\mathbb{N}}
\newcommand{\Z}{\mathbb{Z}}
\newcommand{\Q}{\mathbb{Q}}
\newcommand{\R}{\mathbb{R}}
\newcommand{\C}{\mathbb{C}}
\newcommand{\F}{\mathbb{F}}
\newcommand{\A}{\mathbb{A}}
\newcommand{\PP}{\mathbb{P}}
\newcommand{\E}{\mathbb{E}}
\newcommand{\eps}{\varepsilon}
\newcommand{\dd}{\,\mathrm{d}}
\begin{document}
\section*{UnsolvedMath ID: 30000584}
\textbf{Source problem code:} catalogue entry
\subsection{Definitions and mathematical statement}
In the customary complex-curve setting, let $B$ be a smooth connected projective complex curve of genus $b\geq2$. For $\sigma\in\operatorname{Aut}(B)$ let
\[
\Gamma_\sigma=\{(x,\sigma(x)):x\in B\}\subset B\times B.
\]
A finite subset $S\subset\operatorname{Aut}(B)$ is good if the reduced divisor $\bigcup_{\sigma\in S}\Gamma_\sigma$ is smooth. Since each graph is smooth, this is equivalent to pairwise disjointness of the graphs, or explicitly
\[
\sigma(x)\ne\tau(x)\quad
(\sigma\ne\tau\text{ in }S,\ x\in B).
\]
Equivalently, $\tau^{-1}\sigma$ has no fixed point whenever $\sigma\ne\tau$ are in $S$. The set $S$ is not assumed to be a subgroup.

Its packing number is
\[
\alpha(B,S)=\frac{|S|}{b-1}.
\]
The question asks for its maximal possible value. For a fixed curve the corresponding extremum is
\[
\alpha(B)=\frac{1}{b-1}\max\{|S|:S\subseteq\operatorname{Aut}(B)\text{ good}\}.
\]
For the global interpretation, determine
\[
A=\sup_{\substack{B\text{ smooth projective complex curve}\\g(B)\geq2}}\alpha(B)
\]
and whether this supremum is attained; one can likewise ask for the extremum with genus fixed. These different scopes are distinguished because the extract does not explicitly state which curve or genus parameters are allowed to vary. The complex, smooth, complete interpretation is standard context from the report title, not an explicit field hypothesis in the extract.
\subsection{Short English statement}
How large can the number of pairwise disjoint automorphism graphs be, divided by genus minus one? Good subsets need not be groups. The fixed-curve and global extrema are distinguished because the source does not specify the optimization scope.
\subsection{Sources}
\begin{itemize}
\item Supplied report: \emph{Komplexe Analysis} (2007), Oberwolfach Reports, DOI 10.4171/owr/2006/40, pp. 2399--2446; proposed year 2006. \url{https://doi.org/10.4171/owr/2006/40}.
\item Local UnsolvedMath registry, permanent ID 30000584, catalogue entry; created 25 August 2026, updated 16 September 2026. The corrected\_verified note describes meaning-preserving editorial cleanup. Its supplied 21 August 2026 literature triage labels the question open and reports no located authoritative resolution of the extracted claim. No additional external reference is supplied; the dated assessment is not independently updated here.
\end{itemize}
\end{document}
